\clearpage
\pagestyle{empty}
{\centering
\thispagestyle{empty}
\linespread{1.15}\fontsize{16pt}{22pt}\selectfont\textbf{HỆ THỐNG CHẤM ĐIỂM TÍNH TOÀN VẸN TRÍCH DẪN VÀ PHÁT HIỆN ẢO GIÁC NGUỒN TÀI LIỆU BẰNG BỘ KIỂM TRA LOGIC NƠ-RON -- KÝ HIỆU}\par}

\vspace{0.35\baselineskip}

{\centering\linespread{1.15}\fontsize{16pt}{22pt}\selectfont\textbf{TÓM TẮT}\par}

\vspace{0.45\baselineskip}
\thispagestyle{empty}

Các mô hình ngôn ngữ tạo sinh giúp tăng tốc soạn thảo học thuật nhưng thường sinh ra trích dẫn đúng cú pháp mà trỏ tới tài liệu không tồn tại, bài báo đã rút lại hoặc sai lệch siêu dữ liệu. Tuy nhiên, các công cụ kiểm tra đạo văn và văn bản máy sinh chỉ xét bề mặt câu chữ mà bỏ qua cơ sở dữ liệu đối chứng, trong khi các mô hình học máy thuần túy thất bại trước vi phạm nguồn gốc rời rạc và nhầm lỗi nghẽn mạng với tài liệu bịa đặt. Để lấp khoảng trống này, đồ án tốt nghiệp đề xuất \textbf{CITRA}, nền tảng kiểm chứng trích dẫn theo kiến trúc lai Nơ-ron -- Ký hiệu, thực hiện liên kết hai chiều trên 12 chuẩn trích dẫn và truy hồi liên hợp qua kho cục bộ 129.137 bài báo cùng 15 cơ sở dữ liệu quốc tế. Hệ thống dung hợp đối sánh ngữ nghĩa song ngữ, cây luật ký hiệu phân cấp ưu tiên, hồi quy logistic đa lớp và hiệu chỉnh xác suất Platt kèm cơ chế từ chối phán quyết có chọn lọc. Chúng tôi đánh giá CITRA trên tập kiểm thử độc lập 18.000 trích dẫn (thuộc bộ chuẩn 120.000 mẫu), bộ kiểm thử ngoài phân phối 2.400 mẫu và 310 bản thảo học thuật đa định dạng. Kết quả cho thấy CITRA đạt độ chính xác 99,21\% và Macro-F1 99,04\%, giảm 64,5\% tổng lỗi phân loại so với mô hình học máy cơ sở mạnh nhất ($p < 10^{-44}$) mà không kết tội nhầm bất kỳ trường hợp nào trên 6.480 trích dẫn hợp lệ. Đồng thời, hiệu chỉnh Platt giảm 83,5\% sai số hiệu chỉnh kỳ vọng xuống 0,0411 và các rào chắn ký hiệu nâng độ chính xác ngoài phân phối từ 57,37\% lên 99,96\%. Kết quả nghiên cứu khẳng định việc kết hợp luật nguồn gốc ký hiệu với đối sánh ngữ nghĩa có hiệu chỉnh xác suất mang lại giải pháp kiểm chứng trích dẫn tin cậy và minh bạch cho công tác đánh giá học thuật đại học.

\vspace{0.4\baselineskip}
\noindent\textbf{Từ khóa:} Kiểm chứng tính toàn vẹn trích dẫn, phát hiện ảo giác nguồn tài liệu, trí tuệ nhân tạo Nơ-ron -- Ký hiệu, từ chối phân loại có chọn lọc, hiệu chỉnh xác suất Platt, kỹ thuật phần mềm thực nghiệm.
\clearpage
